% diagram_fcs.tex
% Compact block diagram of the proposed Fuzzy Control System (FCS) for the
% PSO inertia weight w. Designed to be \input{} from paper_wea2026.tex.
%
% Requires: \usepackage{tikz}, \usetikzlibrary{arrows.meta, positioning}
%
% All acronyms used here (PSO, FIS, FCS, PSO-FCS, PSO-STD, COG) MUST be
% defined at first occurrence in the main text BEFORE this figure is included.
%
% Caption is intentionally short: the full reading guide of the diagram
% appears in the paragraph that introduces this figure in paper_wea2026.tex.
%
% diagram_fcs.tex
% Compact LNCS-oriented block diagram of the proposed PSO-FCS architecture.
%


\begin{figure}[H]
\centering
\resizebox{\linewidth}{!}{%
\begin{tikzpicture}[
    x=1cm,y=1cm,
    >=Latex,
    font=\scriptsize,
    block/.style={
        draw, thick, rounded corners=2pt,
        minimum height=0.62cm,
        minimum width=1.85cm,
        align=center,
        inner sep=2pt
    },
    inputblock/.style={
        draw, thick, rounded corners=2pt,
        minimum height=0.52cm,
        minimum width=2.05cm,
        align=center,
        inner sep=1.5pt
    },
    fis/.style={
        draw, thick, rounded corners=2pt,
        minimum height=2.0cm,
        minimum width=3.85cm,
        align=center,
        inner sep=4pt
    },
    sig/.style={->, thick},
    ann/.style={font=\tiny}
]

% --- Fixed coordinates to avoid overlaps after resizing ---
\node[block, minimum width=2.15cm] (pso) at (0,0)
{PSO swarm\\(Eq.~\ref{eq:velocity})};

\node[inputblock] (div) at (2.65,0.48)
{diversity $d^{t}$};

\node[inputblock] (prog) at (2.65,-0.48)
{progress $p^{t}$};

\node[fis] (fis) at (6.15,0)
{\textbf{Mamdani FIS}\\[3pt]
1.~Fuzzification\\[1pt]
2.~\underline{Rule Base}: Rules R1--R8\\[1pt]
3.~Inference Engine\\[1pt]
4.~Aggregation\\[1pt]
5.~Defuzzification};

\node[block, minimum width=2.0cm] (resc) at (9.55,0)
{rescale\\to $[w_{\min},w_{\max}]$};

\node[block, minimum width=1.25cm] (w) at (11.45,0)
{$w^{t+1}$};

% --- Input extraction from PSO ---
\coordinate (tap) at (1.35,0);

\draw[sig] (pso.east) -- (tap) |- (div.west);
\draw[sig] (tap) |- (prog.west);

% --- Inputs to FIS ---
\draw[sig] (div.east) -- ([yshift=0.48cm]fis.west);
\draw[sig] (prog.east) -- ([yshift=-0.48cm]fis.west);

% --- Main output path ---
\draw[sig] (fis.east) -- (resc.west);
\draw[sig] (resc.east) -- (w.west);

% --- Feedback path ---
\draw[sig] (w.south) -- ++(0,-1.00) -| (pso.south)
    node[pos=0.32, below, ann] {feedback to Eq.~\ref{eq:velocity}};

% --- Compact internal specification label ---
%\node[ann, anchor=north] at ($(fis.south)+(0,-0.05)$)
%{rules R1--R8};

\end{tikzpicture}%
}
\caption{PSO-FCS architecture.}
\label{fig:fcs}
\end{figure}